\ExerciseSection

\begin{enumerate}
\def\labelenumi{\arabic{enumi})}
\item
  设 \(f\) 是 \(n\) 个复变量的多项式，系数为复数且不恒等于零。证明 \(\mathbf{C}^n\) 中集 \(S\) 的余集是连通的，这里 S 是由点 \(z = (z_i)\) （满足 \(f(z_1, z_2, \ldots, z_n) = 0\) ；即方程为 \(f = 0\) 的“代数簇”）所组成的集。（若 \(a, b\) 是 \(\mathbb{C}S\) 的两个不同点，考察 \(S\) 与过 \(a\) 和 \(b\) 的复直线的交。）
\item
  设 K 是一个非离散的豪斯多夫拓扑除环（或域）。
\end{enumerate}

\begin{enumerate}
\def\labelenumi{\alph{enumi})}
\tightlist
\item
  设 E 是 K 上的一个 n 维左向量空间。若 \((a_{i})_{1\leqslant i\leqslant n}\) 是 E 的一个基，并且借助双射线性映射把 \(K^{n}\) 的拓扑（诸因子拓扑的乘积）转移到 E 上
\end{enumerate}

\[
(x _ {i}) \rightarrow \sum_ {i = 1} ^ {n} x _ {i} a _ {i},
\]

则这样在 E 上定义的拓扑与所选取的基 \((a_{i})\) 无关，它与 E 的加法群结构相容，并且使映射 \((t, x) \to tx\) （它把 \(K \times E\) 映入 E）连续。若 F 是 E 的一个向量子空间，则 E 的拓扑在 F 上诱导的拓扑与从 F 的任意一个基出发按上述方式在 F 上定义的拓扑相同；F 在 E 中是闭的，并且除非 F = E，CF 在 E 中稠密。

\begin{enumerate}
\def\labelenumi{\alph{enumi})}
\setcounter{enumi}{1}
\tightlist
\item
  把第六章 § 1 的命题 4、5 与 6 推广到除环 K。（为证明当 K 非交换时映射 \(X \to X^{-1}\) 在 K 上每个 n 阶非奇异方阵的邻域内连续，对 n 用归纳法：考察 n 阶单位矩阵 \(I_{n}\) 的一个邻域，注意每个充分接近 \(I_{n}\) 的矩阵 X 都可表成形式
\end{enumerate}

\[
\left( \begin{array}{c c c c} \mathrm{I} & \mathrm{O} & \ldots & \mathrm{O} \\ \lambda_ {2} & & & \\ \vdots & & I _ {n - 1} & \\ \lambda_ {n} & & & \end{array} \right) \left( \begin{array}{c c c c} \mu_ {1} & \mu_ {2} & \ldots & \mu_ {n} \\ \mathrm{o} & & & \\ \vdots & & Y & \\ \mathrm{o} & & & \end{array} \right),
\]

其中 \(I_{n-1}\) 是 n−i 阶单位矩阵，而 Y 是 n−i 阶非奇异矩阵。）

\begin{enumerate}
\def\labelenumi{\alph{enumi})}
\setcounter{enumi}{2}
\item
  若 K 交换且 E 是 K 上秩为 n 的一个代数，则 E 的拓扑与其环结构相容。此外，若 E 有单位元，则 E 的单位群 G 在 E 中是开的且稠密，并且 E 的拓扑在 G 上诱导的拓扑与 G 的群结构相容。（若 e 是 E 的单位元且 x 是 E 的一个单位，则方程 xy = e 与 yx = e 各有且仅有一个解 y；对这两个方程中的每一个，考察给出 y 关于 E 的某个基的分量的等价 n 元线性方程组。）
\item
  若 K 是不完备的域而 E 是 K 上秩为 n 的代数，且 K 的完备化 \(\hat{K}\) 是一个域，则代数 E 的完备化 \(\hat{E}\) 就是把 E 的算子环扩张到 \(\hat{K}\) 而得到的那个代数。由此构造一些完备化不是除环的拓扑除环的例子。
\end{enumerate}

\begin{enumerate}
\def\labelenumi{\arabic{enumi})}
\setcounter{enumi}{2}
\item
  利用第一章 § 3 第 6 号的命题 10，证明实射影空间 \(\mathbf{P}_{n}(\mathbf{R})\) 同胚于 \(\mathbf{P}_{n}(\mathbf{C})\) 的由至少具有一组实齐次坐标的那些点组成的子空间。\([\mathbf{R}_{n+1}^{*}\) 看作嵌入在 \(\mathbf{C}_{n+1}^{*}\) 中，设 A 是 \(\mathbf{R}_{n+1}^{*}\) 中的一个闭集，关于 \(\Delta_{n}(\mathbf{R})\) 是饱和的；设 B 是所有点 \(\zeta x\) 所成的集，其中 \(x \in A\) 且 \(\zeta \in U\) ；证明 \(\overline{\mathbf{B}}\) 关于 \(\Delta_{n}(\mathbf{C})\) 是饱和的，且 A 是 \(\overline{\mathbf{B}}\) 在 \(\mathbf{R}_{n+1}^{*}\) 上的迹。{]}
\item
  在 \(\mathbf{P}_{n}(\mathbf{C})\) 中，设 \(H_{n}\) 是由方程
\end{enumerate}

\[
x _ {0} ^ {2} + x _ {1} ^ {2} + \dots + x _ {n} ^ {2} = 0.
\]

所定义的“二次超曲面”。证明 \(H_{n}\) 的每个点都有一个同胚于 \(C^{n-1}\) 的开邻域，\(H_{n}\) 是连通的，并且 \(H_{n}\) 与任一复射影超平面的余集之交是连通的（为证最后这一断言，化归到 \(n = 2\) 的情形）。

证明 \(H_{2}\) 同胚于 \(S_{2}\) ，而 \(H_{3}\) 同胚于 \(S_{2} \times S_{2}\) （利用二次曲面借助于其直母线的参数表示）。

\begin{enumerate}
\def\labelenumi{\arabic{enumi})}
\setcounter{enumi}{4}
\tightlist
\item
  把球面 \(\mathbf{S}_5\) 看作 \(\mathbf{C}^3\) 的子空间，考察 \(\mathbf{S}_5\) 到 \(\mathbf{R}^7\) 中的映射，它把每个点 \(x = (x_1, x_2, x_3)\) （属于 \(\mathbf{S}_5\) ， \(x_1, x_2, x_3\) 为复数）映到点 \(y = (y_i)_{1 \leqslant i \leqslant 7}\) （属于 \(\mathbf{R}^7\) ），其中
\end{enumerate}

\[
\begin{array}{c} y _ {1} = | x _ {1} | ^ {2} - | x _ {2} | ^ {2}, \quad y _ {2} = \mathcal {R} (x _ {1} \overline {{x}} _ {2}), \quad y _ {3} = \mathcal {I} (x _ {1} \overline {{x}} _ {2}), \quad y _ {4} = \mathcal {R} (x _ {1} \overline {{x}} _ {3}), \\ y _ {5} = \mathcal {I} (x _ {1} \overline {{x}} _ {3}), \quad y _ {6} = \mathcal {R} (x _ {2} \overline {{x}} _ {3}), \quad y _ {7} = \mathcal {I} (x _ {2} \overline {{x}} _ {3}). \end{array}
\]

这个函数在两点 \(\pmb{x}, \pmb{x}'\) （属于 \(\mathbf{S}_5\) ，满足 \(\pmb{x}' = \zeta \pmb{x}\) ，其中 \(|\zeta| = 1\) ）处取相同的值。证明：通过取商，它诱导出 \(\mathbf{P}_2(\mathbf{C})\) 到 \(\mathbf{R}^7\) 的一个子空间上的一个同胚。

¶ 6) 设 K 是一个非离散的豪斯多夫拓扑除环（或域）。

\begin{enumerate}
\def\labelenumi{\alph{enumi})}
\tightlist
\item
  给左射影空间 \(\mathbf{P}_n(\mathbf{K})\) 赋以 \(\mathbf{K}_{n+1}^*\) 的拓扑关于等价关系 \(\Delta_n(\mathbf{K})\) 的商拓扑。把第六章 § 3 的命题 1、4 与 5 推广到赋有此拓扑的 \(\mathbf{P}_n(\mathbf{K})\) 。证明：\(\mathbf{P}_n(\mathbf{K})\) 当 \(\mathbf{K}\) 连通时是连通的，否则是全不连通的。
\end{enumerate}

* b) 若 K 是局部紧的（但非离散的），证明 \(\mathbf{P}_{n}(\mathbf{K})\) 是紧的。（设 U 是 K 中 o 的一个紧邻域；设 \(a \in K\) 使得

\[
\lim _ {m \rightarrow \infty} a ^ {m} = 0;
\]

设 S 是 \(K_{n+1}^{*}\) 中由这样的点 \(\boldsymbol{x} = (x_{i})\) 组成的子集：其所有坐标 \(x_{i}\) 都属于 U，且对某个指标 k 有 \(x_{k} \in a\) . \(\overline{U}\) ；证明 \(\mathbf{P}_{n}(\mathbf{K})\) 是 S 的典范像）。*

\begin{enumerate}
\def\labelenumi{\alph{enumi})}
\setcounter{enumi}{2}
\item
  同样，把第六章 § 3 的定义 2 与命题 6 和 8 推广到 \(\mathbf{P}_{n,p}(\mathbf{K})\) （由 \(p\) 维射影线性簇所成的集，这些簇位于 \(\mathbf{P}_n(\mathbf{K})\) 中）；* 在假设 \(\mathbf{K}\) 局部紧时，也推广命题 7{[}对命题 7，对每个指标列 \(\sigma\) ，考察子集 \(\mathrm{S}_{\sigma}\) —— \(\mathrm{A}_{\sigma}\) 中由所有行都属于 \(b\) ) 中定义的集 S 的矩阵所组成的子集；证明 \(\mathbf{P}_{n,p}(\mathbf{K})\) 是诸集 \(\mathrm{S}_{\sigma}]\) 之并的典范像。* 当 \(\mathbf{K}\) 是域时，推广第六章 § 3 的命题 9。
\item
  若 K 非交换，用 \(\mathbf{P}_{n,p}^{\prime}(\mathbf{K})\) 表示 K 上 n 维右射影空间中 p 维射影线性簇所构成的空间{[}\(\mathbf{P}_{n,p}^{\prime}(\mathbf{K})\) 的拓扑按与 \(\mathbf{P}_{n,p}(\mathbf{K})\) 的拓扑相同的方式定义{]}。证明 \(\mathbf{P}_{n,p}(\mathbf{K})\) 与 \(\mathbf{P}_{n,n-p-1}^{\prime}(\mathbf{K})\) 同胚{[}使每个 \((p+1)\) 维向量子空间 V（左向量空间 \(E=K_{s}^{n+1}\) 的子空间）对应于 V 的正交补，它是 E 的对偶 \(E^{*}\) 的一个向量子空间{]}。
\end{enumerate}

\begin{enumerate}
\def\labelenumi{\arabic{enumi})}
\setcounter{enumi}{6}
\item
  把第六章 § 3 的习题 6 推广到 C 和 H 上的射影线性簇空间。
\item
  把第六章 § 3 的习题 7、8 与 9 推广到空间 \(\mathbf{W}_{n,p}\) （第 4 号）。
\end{enumerate}

在四元数体上的向量空间 \(H^{n}\) 中，我们仍用 \(W_{n,p}\) 表示所有序列 \((x_{k})_{1\leqslant k\leqslant p}\) （由 p 个向量 \(x_{k}=(x_{kj})_{1\leqslant j\leqslant n}\) 组成）的全体所成的集，这些向量满足

\[
\sum_ {j = 1} ^ {n} x _ {i j} \overline {{{x}}} _ {i j} = \mathrm{I} \quad (\mathrm{I} \leqslant i \leqslant p)
\]

且

\[
\sum_ {j = 1} ^ {n} x _ {i j} \bar {x} _ {i j} = 0 \quad (i \neq k)
\]

（ \(\bar{x}\) 表示四元数 \(x\) 的共轭）。把第六章 § 3 的习题 7、8 与 9 推广到这些空间 \(\mathbf{W}_{n,p}\) 。

